\subsection{Introducing Sets} 
Sets are one of the fundamental concepts of mathematics, but it turns out to be remarkably difficult to define them precisely. Here's an informal definition:
\begin{definition}[Set (informal)]
A \term{set} is a collection of (usually mathematical) objects. The objects in a set are called its \term{elements}.
\end{definition}
This definition is not helpful. Let's see some examples.
\begin{example}
We can have a set containing the numbers 1, 2, and 3. We would write this as $\{1,2,3\}$. We could also give it a name, for example, $A=\{1,2,3\}$.
\end{example}
Key things to notice about this:
\begin{itemize}
\item There are \makeblanksmall numbers in this set: \makeblanksmall, \makeblanksmall, and \makeblanksmall.
\item However, this is \makeblanksmall set, named \makeblanksmall.
\item Think of the set as a ``container'' full of numbers.
\item The number $2$ is an \makeblank of the set \makeblanksmall (it's one of the numbers on the list). To indicate this, we can write $2\in A$.
\item The number $5$ is \emph{not} an \makeblank of the set \makeblanksmall (because it's not one of the numbers on the list). To indicate this, we can write $5\notin A$.
\end{itemize}
When we talk about sets, \makeblank doesn't matter. We could list this set as any of:
\begin{align*}
A&=\{1,2,3\}&
A&=\{1,3,2\}&
A&=\{2,1,3\}\\
A&=\{2,3,1\}&
A&=\{3,1,2\}&
A&=\{3,2,1\}
\end{align*}
It's also the case that \makeblank doesn't matter. All of these are also the same set:
\begin{align*}
A&=\{1,2,3\}&
A&=\{1,1,2,3\}&
A&=\{1,2,2,3,3,3,1\}
\end{align*}
All that matters is whether a number \emph{is} or \emph{is not} in the set.